\section{Discussion}
\label{sec:discussion}

\subsection{Interpreting the effect on cJSON}
The result in Section~\ref{sec:evaluation} says something narrower than ``the
LLM makes fuzzing better'': across fifteen seeded runs, refinement reliably
moved a generator from producing mostly-invalid input toward producing
mostly-valid input, using only the three coverage-free proxies it was given.
That acceptance-rate effect is real and statistically clean at the sample
size tested -- the scale-free metric on which the two arms are directly
comparable -- but it is worth being explicit about what it does not show.
A 97.1\% mean acceptance rate is close to a ceiling, and a
generator that always emits valid input has, by construction, stopped
exercising the accept/reject boundary at all -- exactly the region the
project's own \texttt{near\_valid\_json} baseline strategy exists to probe.
Acceptance rate rising is desirable when the baseline is far from valid and
mostly testing the rejection path, which is the regime our baseline started
in; it is not obviously desirable in the limit, where a generator that never
produces a rejected or malformed input has traded away exactly the
boundary-probing behavior that historically finds parser bugs. We did not
observe this failure mode directly -- within the refined arm, structural
fingerprint counts also rose and rejection signatures did not collapse to a
single trivial case, though these counts are descriptive rather than
tested against the baseline (Section~\ref{sec:evaluation}) -- but the proxy
signal as designed rewards acceptance rate monotonically and has no
explicit floor protecting the near-valid regime.
A refinement objective that also rewarded a target reject rate, rather than
treating ``more accepted'' as strictly better without limit, would be a
direct, low-cost improvement to the design evaluated here.

\subsection{Interpreting the null result on parson}
Finding no crash across 3{,}000 sanitizer-instrumented executions against
parson is evidence about this pipeline's search within this budget on this
pinned commit, not evidence about parson's correctness in general, and not
evidence about the method's bug-finding power in general. We treated the
source-level argument in Section~\ref{sec:evaluation} as a way to make the
absence explainable rather than mysterious, but an explainable absence is
still an absence: five iterations against two hand-picked stress axes (wide
structures, and the internal free path) is a small, targeted search of a
small, well-tested codebase, and the honest reading is that this budget was
not enough to distinguish ``no bug in the code paths we exercised'' from ``a
bug exists in a code path we did not exercise,'' such as parson's public
mutation API, which the black-box parse-then-free harness used here cannot
reach at all.

What the parson run does establish is narrower and more mechanical: every
non-target-specific pipeline component (\texttt{runner}, \texttt{campaign},
\texttt{proposal}, \texttt{refinement}, \texttt{triage}, \texttt{minimize})
worked unmodified against a second parser with a different, and in two
respects opposite, set of grammar adaptations, with only a new harness and
build script required. That is a claim about the pipeline's portability
across targets, not a claim about its effectiveness at finding bugs in any
particular one of them.

\subsection{Threats to validity}
\textbf{Construct validity.} A ``crash'' in this pipeline is defined purely
by sanitizer output and process signals; it cannot detect a parser that
silently misparses, accepts input the grammar disallows, or otherwise
violates the specification without a memory-safety fault. The grammar
mismatches documented in Sections~\ref{sec:methodology}
and~\ref{sec:evaluation} (duplicate keys, full-buffer consumption) were
found by deliberately generating inputs in that region and reading the
result, not by any automated oracle; a generator that never targets an
under-specified region of the grammar would not surface this class of
finding at all.

\textbf{Internal validity.} The structural-fingerprint proxy truncates at
256 bytes (Section~\ref{sec:methodology}), so it systematically undercounts
diversity among longer, deeply nested inputs -- exactly the inputs a
refinement loop chasing ``more diversity'' is most likely to produce over
successive iterations, which could understate the true diversity signal
reflected in the descriptive fingerprint counts we report.

\textbf{Statistical validity.} $n{=}15$ runs per arm is still modest; the
exact test used is appropriate for that size and makes no distributional
assumption, but $p \approx 1.29\times10^{-8}$ on the acceptance-rate
comparison is a ceiling on how small an exact $p$-value fifteen-per-arm
observations can produce, not a measure of effect size on its own, and we report the actual
acceptance-rate numbers, alongside the descriptive fingerprint and
rejection-signature counts, in Table~\ref{tab:cjson-comparison} for that
reason. Baseline and refined arms also executed different total numbers of
inputs (Section~\ref{sec:evaluation}), which is why only the scale-free
acceptance rate is tested statistically; any reader extending this
comparison with new count-based metrics should convert to a rate first.

\textbf{External validity.} All results are on JSON, on two relatively
small and independently well-tested C parsers, with one LLM
(\texttt{gpt-4.1-mini}) at one temperature setting for the only
LLM-driven arm. We make no claim that the magnitude of the cJSON effect, or
the absence of a crash on parson, generalizes to other formats, other
parsers, or other proposer models.

\subsection{When this approach is, and is not, the right tool}
The motivating case for a coverage-free feedback loop is a target that
genuinely cannot be instrumented -- a vendored dependency pinned at a
specific commit for compatibility reasons, a license-restricted or
closed-source binary, or a cross-compiled/embedded parser where an
instrumented build would not be representative of what ships. Where
coverage instrumentation \emph{is} available, a coverage-guided or
grammar-aware coverage-guided fuzzer (Section~\ref{sec:related-work}) has
strictly more signal to work with and should be preferred; this pipeline is
not proposed as a replacement for that setting, only as an option for the
black-box one.

\subsection{Future work}
Several extensions are direct given the existing infrastructure. First, the
feedback-signal ablation in Section~\ref{sec:evaluation-ablation} was run at
a small scale (four usable trials per mode) and its rejection-signature
finding is reported descriptively for that reason; a larger-sample
replication is the natural next step to establish whether the pattern we
observed holds statistically. Second, adding an explicit reject-rate floor
to the refinement prompt, motivated directly by this section's concern about
an acceptance-rate ceiling and now further motivated by the ablation's
rejection-signature pattern, could test whether it preserves
boundary-probing behavior without sacrificing the acceptance-rate gain
measured here. Third, and most directly relevant to the parson null result, driving
\texttt{json\_value\_free} through parson's public mutation API rather than
only the parse-then-free path would exercise object lifetime code this
black-box harness cannot currently reach at all. Fourth, comparing against
an existing external fuzzer under an equalized time or execution budget,
rather than only against this project's own static baseline, is the most
direct way to situate the magnitude of the cJSON result against prior work.

Two further extensions are larger in scope and deliberately left for future
work rather than attempted here. One is broadening beyond a second
same-language, same-format target: a third C JSON parser would only weakly
extend the generalization claim already established by parson, whereas a
target in a different format, or a memory-safe language where a
``crash'' oracle must be redefined around uncaught exceptions and API
contract violations rather than sanitizer output -- for instance, JSON
deserialization in a Dart/Flutter mobile application -- would test
generalization along a materially different axis than implementation alone.
The other is applying this pipeline to a live, running mobile or server
application rather than a bounded harness process, which would require
redesigning the outcome-classification step in
Section~\ref{sec:methodology} around that environment's own failure
signals.
